Large language model (LLM) agents can draft analog circuits, but their evaluations rarely test what makes analog
design hard: meeting many quantitative specifications at once, in a real process, across process, voltage and
temperature (PVT) corners. We present SPECTRA, an agent that turns a natural-language specification into a sized,
PVT-verified SKY130 netlist. Instead of asking the LLM for transistor widths, SPECTRA has it write a
\emph{design-intent program} that sizes every device from $(g_m/I_D, L, I_D)$ through PDK lookup tables and predicts
its own performance. A deterministic layer executes and simulates the program and returns a device-level diagnosis;
the LLM uses it to choose between rewriting the program and handing it to a numerical refiner that searches the
program's low-dimensional intent space, with adaptive worst-condition selection for robust tasks. We also release
SPECTRA-Bench: \NTasksTotal{} specification-driven tasks over twelve circuit families, with certified specifications,
an independently written designer tier, and verification over 45 PVT conditions. On a preregistered external track
of \SizTasks{} fixed-topology tasks, seven paper-backed framework adapters and three seeds share a 300-simulation,
12-call budget. SPECTRA reaches TT/PVT success \SizSpectraTT{}/\SizSpectraPVT{} versus
AnaFlow's \SizAnaFlowTT{}/\SizAnaFlowPVT{} in that frozen matrix. We therefore report a benchmark-scoped
comparison, together with cross-simulator and capability audits, rather than a universal SOTA claim. The diagnosis
does not by itself improve fault localisation, and intent programs mispredict second-order effects such as supply
rejection.
